\documentclass[10pt,a4paper]{amsart}
\usepackage[margin=19mm]{geometry}
\usepackage{amsmath,amssymb,mathtools}
\usepackage[scr=rsfs,cal=euler]{mathalfa}
\usepackage{xfrac}
\usepackage[unicode,hidelinks,bookmarks=false]{hyperref}
\newcommand{\ie}{i.e.\ }
\newcommand{\cf}{cf.\ }
\newcommand{\resp}{respectively\ }
\newcommand{\Subsection}{\subsection}
\providecommand{\nobreakdash}{\nobreak\hbox{-}}
\setlength{\emergencystretch}{2em}
\multlinegap0pt
\usepackage{amssymb}
\newtheorem{theorem}{Theorem}
\newtheorem{lemma}{Lemma}
\newcommand{\B}{\mathbb{B}}
\newcommand{\C}{{\mathbb{C}}}
\renewcommand{\P}{{\mathbb{P}}}
\title{Meromorphic mappings on K\"{a}hler manifolds weakly sharing hyperplanes in $\P^n(\C)$}
\author{Si Duc Quang}
\date{Source: arXiv:2405.06268v1 (CC BY 4.0)}
\begin{document}
\maketitle

\noindent\emph{Source and licence.} Meromorphic mappings on K\"{a}hler manifolds weakly sharing hyperplanes in $\P^n(\C)$. Version arXiv:2405.06268v1 (CC BY 4.0), distributed by arXiv under the Creative Commons Attribution 4.0 International licence, http://creativecommons.org/licenses/by/4.0/, as recorded in the arXiv licence metadata for this exact version and shown on the abstract page https://arxiv.org/abs/2405.06268v1. Full licence notice: \texttt{licenses/arxiv-2405.06268-CC-BY-4.0.txt}.\par\smallskip

\noindent\textit{Editorial scope.} Counted unit: the article's theorem 1.1 (source0.tex lines 121--130). The article's complete original proof of it is delivered with it (source0.tex lines 366--400, one \texttt{proof} environment of 35 lines, which the article places away from the statement). No mathematics is written, repaired or re-derived: the statement and the proof are the article's own lines, in the article's own notation, and no retained line is substituted or re-flowed.  Retained uncounted as the source-local context the proof needs: lines 318--332.  Nothing was omitted from any retained span, so the package carries no omission record.  No retained line contains a citation, so the wrapper carries no bibliography and no bibliography entry.  No retained line contains a figure environment, an image-inclusion command or drawing code, so no image is delivered and the package declares no asset dependency.  Editorial formatting only, in addition: an AMS \texttt{amsart} preamble that restates the article's own declarations and macros used by the retained lines, namely the environments \texttt{theorem}, \texttt{lemma} on separate counters, together with the standard packages those lines need.  The retained environments print in this document as Theorem 1, Lemma 1 in order of appearance, whereas the article prints the same items with the numbering of its own full document.  The complete native source of the article as distributed in the arXiv e-print for this exact version is bundled unchanged as \texttt{sources/158946/source0.tex}.
\begin{lemma}\label{3.4}
Let $M$ be a complete, connected K\"{a}hler manifold whose universal covering is biholomorphic to $\B(R_0) \subset \C^m$, where $0<R_{0} \leq \infty$. Let $f, g: M \rightarrow\P^n(\C)$ be linearly nondegenerate meromorphic mappings, which satisfy the condition $(C_{\rho_f}), (C_{\rho_g})$ for non-negative constants $\rho_f,\rho_g$, respectively. Let $H_1,\ldots,H_q$ be $q$ hyperplane of $\P^n(\C)$ in general position and let $n+2\le p\le 2n+2<q$. Assume that:
\begin{itemize}
\item[(a)] $\dim f^{-1}(H_i\cap H_j)\le m-2 \ \forall\ 1\leq i< j\leq  q,$
\item[(b)] $f^{-1}(H_i)=g^{-1}(H_i)\ \forall\ 1\le i\le p,\ f^{-1}(H_i)\subset g^{-1}(H_i)\ \forall\ p+1\le i\le q,$
\item[(c)] $f=g$ on $\bigcup_{i=1}^qf^{-1}(H_i).$
\end{itemize}
Then $f\equiv g$ if there exist non-negative numbers $t\le\frac{2}{n}$ and $\alpha$ such that:
\begin{itemize}
\item[(1)] $(2+t)(q-n-1)-(2n+2+p(t+\alpha))> 0,$
\item[(2)] $\left(2+\frac{\alpha}{n}\right)(p-n-1)-(2n+2)> 0,$
\item[(3)] $\frac{(2+t)(q-n-1)-(2n+2+p(t+\alpha))}{\rho_f}+\frac{\left(2+\frac{\alpha}{n}\right)(p-n-1)-(2n+2)}{\rho_g}> 2\left((2+t)\ell_f+(2+\frac{\alpha}{n})\ell_g\right),$
\end{itemize}
where $t=0$ if $p> q-n-1$ and $t=\frac{2}{n}$ if $p\le q-n-1$.
\end{lemma}

\begin{theorem}\label{1.1}
Let $M$ be a complete, connected K\"{a}hler manifold whose universal covering is biholomorphic to $\B(R_0) \subset \C^m\ (0<R_0\le +\infty)$. Let $f, g: M \rightarrow\P^n(\C)$ be linearly nondegenerate meromorphic mappings satisfying a condition $(C_{\rho})$. Let $H_1,\ldots,H_q$ be hyperplanes of $\P^n(\C)$ in general position with $\dim f^{-1}(H_i\cap H_j)\le m-2$ for every $1\le i<j\le q$, such that
\begin{itemize}
\item[(i)] $f^{-1}(H_i)=g^{-1}(H_i)\ \forall\ 1\le i\le p$, $f^{-1}(H_i)\subset g^{-1}(H_i)\ \forall\ p+1\le i\le q$, 
\item[(ii)] $f=g$ on $\bigcup_{i=1}^qf^{-1}(H_i)$, 
\end{itemize}
where $n+2\le p\le 2n+2$. Then $f\equiv g$ if 
$$q>2n+2+pn\left(\frac{n+1}{p-n-1}-1\right)+2\rho\left(\ell_f+\frac{n+1}{p-n-1}\ell_g\right)$$
$$\text{or }\ q>2n+1+pn\left(\frac{n}{p-n-1}-\frac{n-1}{n}\right)+2\rho\left(\ell_f+\frac{n}{p-n-1}\ell_g\right).$$
\end{theorem}

\begin{proof}[Proof of Theorem \ref{1.1}] By Theorem \ref{3.4}, in order to show that $f\equiv g$ we just show the existence of the non-negative numbers $t\le\frac{2}{n}$ and $\alpha$ satisfying the inequalities (1), (2), (3) of Theorem \ref{3.4}. We choose
$$ \alpha:=n\left(\frac{2n+2+\epsilon}{p-n-1}-2\right),$$
where $\epsilon$ is a positive constant small enough. Then the inequality (2) of Theorem \ref{3.4} is satisfied. The inequalities (1) and (3) of Theorem \ref{3.4} become
\begin{align}\label{4.1}
(2+t)(q-n-1)-2n-2-p\left(t-2n+n\frac{2n+2+\epsilon}{p-n-1}\right)> 0,
\end{align}
and
\begin{align}
\label{4.2}
\begin{split}
(2+t)(q-n-1)-2n-2-p\left(t-2n+n\frac{2n+2+\epsilon}{p-n-1}\right)+\epsilon\\
>2\rho\left((2+t)\ell_f+\frac{2n+2+\epsilon}{p-n-1}\ell_g\right).
\end{split}
\end{align}
Therefore, in order to show the existence of such $t$ and $\alpha$ (equivalent to the existence of $\epsilon >0$) we just to show that there is $t\in[0,\frac{2}{n}]$ such that:
\begin{align}
\label{4.3}
\begin{split}
(2+t)(q-n-1)-2n-2-&p\left(t-2n+n\frac{2n+2}{p-n-1}\right)\\
&>2\rho\left((2+t)\ell_f+\frac{2n+2}{p-n-1}\ell_g\right).
\end{split}
\end{align}
If put $t=0$ then the inequality (\ref{4.3}) is equivalent to that 
\begin{align*}
q-n-1&>n+1+pn\left(\frac{n+1}{p-n-1}-1\right)+2\rho\left(\ell_f+\frac{n+1}{p-n-1}\ell_g\right)\\ 
\Leftrightarrow\ q&>2n+2+pn\left(\frac{n+1}{p-n-1}-1\right)+2\rho\left(\ell_f+\frac{n+1}{p-n-1}\ell_g\right).
\end{align*}
If put $t=\frac{2}{n}$ then the inequality (\ref{4.3}) is equivalent to that 
\begin{align*}
\frac{n+1}{n}(q-n-1)&>n+1+p\left(\frac{1-n^2}{n}+n\frac{n+1}{p-n-1}\right)+2\rho\left(\frac{n+1}{n}\ell_f+\frac{n+1}{p-n-1}\ell_g\right)\\
\Leftrightarrow\ q&>2n+1+pn\left(\frac{n}{p-n-1}-\frac{n-1}{n}\right)+2\rho\left(\ell_f+\frac{n}{p-n-1}\ell_g\right).
\end{align*}
Then, with the asumption of the themrem, the inequality (\ref{4.3}) is satisfied for $t=0$ or $t=\frac{2}{n}$. Hence, $f\equiv g$.
The proof of the theorem is completed.
\end{proof}

\end{document}
